% Propietario conservado: /Users/ruben/Documents/New project/output/REVISION_ARGUMENTAL_APERTURAS_CIERRES_20260915/III/spanish_source/manuscrito/sections/07_carga_y_cuantos.tex
\section{Charge section and electrical quantization relations}
\label{iii:iii:sec:carga}

The constitutive response establishes a relation between the two sheets and a
relative propagation scale. The charge section is obtained by composing
this response with the fine-structure constant and the action, both
constructed in the preceding chapters. This order has a precise
consequence: charge does not retrospectively determine the dodecaphase record.
It is the record, together with the action and the vacuum response, that
determines the charge section on the corresponding line.

\subsection{Positive determination of the charge section}

Consider the positive sections $\hbar_a$, with
$a\in\{\mathrm{pre},\mathrm{ret}\}$, and $h_a=2\pi\hbar_a$.
The impedance $Z_{\rm vac}$ has already been constructed through
\eqref{iii:iii:eq:dimensions}. The quotient $h_a/Z_{\rm vac}$ belongs to
the quadratic charge line. This dimensional compatibility allows
the following definition to be formulated without yet choosing an SI chart.

\begin{definition}[Positive charge section]
The charge section corresponding to action orientation $a$ is
\begin{equation}
 e_a=\left(\frac{2\alpha h_a}{Z_{\rm vac}}\right)^{1/2}.
 \label{iii:iii:eq:charge}
\end{equation}
The root is taken on the positive charge half-line. Its conjugate section
is $-e_a$; the sign of the latter is not identified with the action
return index.
\end{definition}

\begin{proposition}[Existence, uniqueness, and compatibility of the coupling]
The equation $Z_{\rm vac}e_a^2=2\alpha h_a$ determines a unique
positive section. In the constitutive realization, one also has
\begin{equation}
 \alpha=\frac{Z_{\rm vac}e_a^2}{2h_a}
 =\frac{e_a^2}{4\pi\varepsilon_{\rm vac}\hbar_a c_{\rm vac}}.
 \label{iii:iii:eq:coupling}
\end{equation}
\end{proposition}
\begin{proof}
Positivity of $\alpha$, $h_a$, and $Z_{\rm vac}$ guarantees a unique
positive root. The constitutive identities give
$\varepsilon_{\rm vac}c_{\rm vac}=Z_{\rm vac}^{-1}$.
Substituting this equality and $h_a=2\pi\hbar_a$ into the first expression
of \eqref{iii:iii:eq:coupling} yields the second.
\end{proof}

The last expression is the usual form of electromagnetic coupling
in a physical chart. Here it appears as a subsequent identity of the constructed
sections. The proof does not recalculate $\alpha$ from a measured
charge: it establishes the algebraic consistency of the preceding result with the
section just defined. Physical identification additionally requires
retaining the chart of units and the coupling regime used in the
comparison. An identity between quantities and an experimental comparison
have different functions in that recognition.

\subsection{Resistance, conductance, flux, and frequency}

The following four quantities use the same section $e_a$. Their
dependencies can be fully resolved before any decimal is evaluated.

\begin{theorem}[Composition of the electrical quanta]
The sections
\begin{align}
 R_K&=\frac{h_a}{e_a^2}=\frac{Z_{\rm vac}}{2\alpha},
 &G_0&=\frac{2e_a^2}{h_a}=\frac{4\alpha}{Z_{\rm vac}},
 \label{iii:iii:eq:rk-g0}\\
 \Phi_{0,a}&=\frac{h_a}{2e_a}
 =\sqrt{\frac{h_aZ_{\rm vac}}{8\alpha}},
 &K_{J,a}&=\frac{2e_a}{h_a}=\Phi_{0,a}^{-1}
 \label{iii:iii:eq:flux-josephson}
\end{align}
satisfy exactly
\begin{equation}
 R_KG_0=2,\qquad G_0Z_{\rm vac}=4\alpha,\qquad
 K_{J,a}\Phi_{0,a}=1,\qquad
 K_{J,a}^{2}R_K=\frac4{h_a}.
 \label{iii:iii:eq:quanta-identities}
\end{equation}
\end{theorem}
\begin{proof}
Substituting $e_a^2=2\alpha h_a/Z_{\rm vac}$ into the first two
definitions cancels $h_a$ and gives \eqref{iii:iii:eq:rk-g0}. Squaring
the positive section $h_a/(2e_a)$ gives
$h_aZ_{\rm vac}/(8\alpha)$; positivity fixes the root in
\eqref{iii:iii:eq:flux-josephson}. Its inverse is $2e_a/h_a$.
The first three products in \eqref{iii:iii:eq:quanta-identities}
reduce, respectively, to $2$, $4\alpha$, and $1$. Finally,
$(4e_a^2/h_a^2)(h_a/e_a^2)=4/h_a$.
\end{proof}

The physical realization recognizes in these sections the von
Klitzing resistance, the conductance quantum, the flux quantum, and the
Josephson constant. The factor $2$ in the conductance belongs to the normalization
of the quantity defined here; identifying it with a multiplicity of
channels in a material system requires specifying that representation.
The four equations do not constitute four independent generators:
they share $\alpha$, the action, and the vacuum response. Precisely this
dependency produces the cross-identities and makes it possible to check an
evaluation without altering the operations from which it arose.

\subsection{Action return and transport of the sections}

The quotient $R_{\rm act}=h_{\rm pre}/h_{\rm ret}$, obtained before the
angular pair, preserves the distinction between the two action sections.
It does not change the sheet assignments of the constitutive response. These are
two distinct operations: action return and channel interchange.

\begin{proposition}[Covariance under return]
With common $\alpha$ and $Z_{\rm vac}$, one has
\begin{align}
 \frac{e_{\rm pre}}{e_{\rm ret}}&=R_{\rm act}^{1/2},
 &R_K^{\rm pre}&=R_K^{\rm ret},
 &G_0^{\rm pre}&=G_0^{\rm ret},\label{iii:iii:eq:charge-return}\\
 \frac{\Phi_{0,\rm pre}}{\Phi_{0,\rm ret}}&=R_{\rm act}^{1/2},
 &\frac{K_{J,\rm pre}}{K_{J,\rm ret}}&=R_{\rm act}^{-1/2}.
 \label{iii:iii:eq:flux-return}
\end{align}
\end{proposition}
\begin{proof}
The first identity follows by dividing the two versions of
\eqref{iii:iii:eq:charge} and taking the positive root. In
\eqref{iii:iii:eq:rk-g0}, the action has already cancelled. The other two
identities follow from the powers $h_a^{1/2}$ and $h_a^{-1/2}$ in
\eqref{iii:iii:eq:flux-josephson}.
\end{proof}

Resistance and conductance preserve the constitutive quotient, whereas
flux and frequency also retain the action orientation. This
separation gives a structural reading of their dependencies: some quantities
are invariants of the sheet quotient, while others additionally measure the
action section used. The metrological comparison must specify which section it evaluates;
the numerical proximity of the two actions does not authorize their identification.
